\documentclass{report}
\usepackage{amsmath,amssymb}
\setlength{\parindent}{0mm}
\setlength{\parskip}{1em}
\begin{document}
\begin{center}
\rule{6in}{1pt} \
{\large Darko Volkov \\
{\bf An inverse problem for the recovery of active faults from surface observations}}

Department of Mathematical Sceinces \\ WPI \\ 100 Institute road \\ Worcester MA 01609
\\
{\tt darko@wpi.edu}\end{center}

The forward modeling of the quasi-static evolution of faults in the
Earth�s upper crust has been extensively studied by geophysicists. The
use of GPS and InSAR technologies has made it possible to measure surface
displacement fields with remarkable accuracy.
We will review in this talk the mathematical and numerical inverse
problem consisting of portraying faults based on surface measurements. In
the quasi static case we were able to recover the average position and
orientation of faults as well as total slips across these faults.
Reconstructing averages induces regularity in the inverse problem and
leads to robust reconstruction algorithms for the 2D and the 3D model.
In more recent work we consider the case of shallow faults and abundant
(but noisy) surface data. In this case a more complete reconstruction is
possible: the recovered fault is expressed as a finite set of poles. We
are currently investigating the time dependent case based on the
assumption that stability failure propagates in time along the fault.


\end{document}
